\chapter{Quick Start}
\label{chap:quick_start}

This chapter is intended to get a new user running their first simulation
within a few minutes. The example below illustrates a basic
frequency--severity simulation using \actsim's \texttt{StochasticSimulator}.

\section{Minimal Example}

The following script simulates 10{,}000 years of aggregate losses, using a
Poisson frequency distribution with mean $50$ and a lognormal severity
distribution.

\begin{lstlisting}[style=pythonstyle]
from actsim import StochasticSimulator

simulator = StochasticSimulator(
    freq_dist="poisson",
    freq_params=(50,),
    sev_dist="lognormal",
    sev_params=(10, 0.5),
    num_sim=10000,
    keep_all=True,
    seed=1234,
)

# Run the simulation
simulator.gen_agg_simulations()

# Inspect the simulated aggregate losses
print(simulator.results.head())

# Calculate the 99th percentile of the aggregate loss distribution
var_99 = simulator.calc_agg_percentile(99)
print(f"99% VaR: {var_99:,.2f}")
\end{lstlisting}

\section{What the Example Does}

The example performs the following steps:

\begin{enumerate}[leftmargin=*]
    \item Constructs a \texttt{StochasticSimulator} configured with a
          Poisson frequency distribution and a lognormal severity
          distribution.
    \item Sets the number of simulated years to 10{,}000.
    \item Sets \texttt{keep\_all=True} so that individual event-level data
          are stored alongside the aggregate results, enabling later
          calculation of metrics such as Occurrence Exceedance Probability.
    \item Fixes the random seed to 1234 to make the simulation reproducible.
    \item Calls \texttt{gen\_agg\_simulations()} to run the simulation.
    \item Calculates the 99th percentile of the aggregate loss
          distribution.
\end{enumerate}

\section{Next Steps}

After running this minimal example, users may wish to:

\begin{itemize}[leftmargin=*]
    \item Explore the \texttt{DistributionFitter} class for fitting
          distributions to observed loss data
          (Chapter~\ref{chap:distribution_fitter}).
    \item Use the \texttt{ClaimSimulator} class to generate synthetic
          claim-level data from a policy table
          (Chapter~\ref{chap:claim_simulator}).
    \item Apply deductibles, limits, and other contract terms to the
          simulated losses (Section~\ref{sec:deductibles_and_limits}).
    \item Integrate the simulated data with the \texttt{chainladder} package
          for reserving analysis (Chapter~\ref{chap:chainladder}).
\end{itemize}

The remainder of the manual covers each capability in detail, with worked
examples and validation guidance.
